\section{Methodology}
\label{methodology}

\subsection{Framework Overview}
\label{sec:method-overview}

The proposed method transforms a natural language requirement into an executable requirement representation.
The workflow contains four stages: requirement semantic analysis, knowledge-guided task transformation, execution-flow construction, and human review.
The output is a structured plan that links each requirement element to tasks, actions, constraints, and dependencies.

\subsection{Multi-Agent Task Planning Workflow}
\label{sec:method-agents}

\subsubsection{Analyst Agent}
\label{sec:method-analyst}

The analyst agent identifies the functional intent, relevant entities, constraints, and observable behaviors expressed in the requirement.
Its output is a semantic summary used by later planning steps.

\subsubsection{Planning Agent}
\label{sec:method-planning-agent}

The planning agent decomposes the analyzed requirement into candidate tasks and orders them according to dependency and constraint information.
It is responsible for constructing an action sequence that can be inspected and executed by an intelligent system.

\subsubsection{Execution Agent}
\label{sec:method-execution-agent}

The execution agent checks whether the task workflow is complete.
It inspects whether the plan contains actionable steps, inputs, outputs, and dependencies.

\subsection{Requirement Semantic Analysis}
\label{sec:method-semantic}

\subsubsection{Functional Intent}
\label{sec:method-intent}

Functional intent captures the capability or behavior expected by the requirement.
It defines the goal that the task workflow must preserve.

\subsubsection{Constraints}
\label{sec:method-constraints}

Constraints capture conditions that restrict task execution, including input conditions, business rules, quality constraints, and stakeholder-specified limitations.

\subsubsection{Action Semantics}
\label{sec:method-action-semantics}

Action semantics identify the operations implied by the requirement, such as collecting information, validating conditions, updating state, notifying users, or producing outputs.

\subsubsection{Execution Dependencies}
\label{sec:method-dependencies}

Execution dependencies describe ordering relations among tasks and actions.
They help prevent plans that satisfy isolated actions but fail as complete workflows.

\subsection{Knowledge-Guided Task Transformation}
\label{sec:method-knowledge}

\subsubsection{Requirement Knowledge}
\label{sec:method-req-knowledge}

Requirement knowledge provides domain concepts, common requirement types, and terminology used to interpret stakeholder statements.

\subsubsection{Task Templates}
\label{sec:method-task-templates}

Task templates define reusable structures for common executable requirement patterns, such as input validation, authorization checks, notification workflows, and state transitions.

\subsubsection{Planning Rules}
\label{sec:method-planning-rules}

Planning rules constrain task decomposition and ordering.
They specify when tasks require preconditions, postconditions, dependency checks, or human confirmation.

\subsubsection{Prompt Strategies}
\label{sec:method-prompts}

Prompt strategies guide the agents to produce structured intermediate artifacts instead of unconstrained natural language summaries.
Prompts are designed around intent extraction, action identification, dependency reasoning, and plan consistency checking.

\subsection{Executable Requirement Construction}
\label{sec:method-construction}

\subsubsection{Task Decomposition}
\label{sec:method-task-decomposition}

The method decomposes each requirement into a set of tasks that are specific enough to be inspected, assigned, or executed by an intelligent agent.

\subsubsection{Action Sequence Generation}
\label{sec:method-action-sequence}

Tasks are ordered into an action sequence that reflects preconditions, data dependencies, and expected outputs.

\subsubsection{Requirement-to-Task Mapping}
\label{sec:method-mapping}

Each task is linked back to the requirement fragment that justifies it.
This mapping supports traceability between stakeholder intent and executable behavior.

\subsubsection{Execution Flow Construction}
\label{sec:method-flow}

The final representation organizes tasks, actions, constraints, and dependencies into an execution flow.
The flow is intended for planning and inspection, not for deployment-level execution.

\subsection{Human-in-the-Loop Interaction}
\label{sec:method-hitl}

Human reviewers inspect ambiguous mappings, missing constraints, and inconsistent task orderings.
Their feedback is used to revise the executable requirement representation before evaluation.
